%%%%%%%%%%%%%%%%%%%%%%%%%%%%%%%%%%%%%%%%%%%%%%%%%%%%%%%%%%%%%%%%%%%%%%
% Lớp 10 | Hàm số bậc hai và đồ thị | Hàm số và đồ thị
% Dạng 0D3?1-3: Giá trị của hàm số
% Số câu: 16  (sinh tự động bởi tach_cau.py; từ nay thầy sửa trực tiếp file này)
%%%%%%%%%%%%%%%%%%%%%%%%%%%%%%%%%%%%%%%%%%%%%%%%%%%%%%%%%%%%%%%%%%%%%%

%%%%%%%%%%%%%%%% TRẮC NGHIỆM 4 PHƯƠNG ÁN %%%%%%%%%%%%%%%%

% nguồn: MaHoaNH_TracNghiem_Lop10
\begin{ex}%[0D3N1-3]%[ma:82fd3d4f]
	Cho hàm số $y=f(x)=2x^2-1$. Tính  $f(2)$.
	 \choice
		{$f(2)=2$} 
		 {$f(2)=3$}
		 {\True $f(2)=7$}
		 {$f(2)=5$}
	 \loigiai{Ta có $f(2)=2\cdot2^2-1=7$.}
\end{ex}

% nguồn: MaHoaNH_TracNghiem_Lop10
\begin{ex}%[0D3N1-3]%[ma:f2803899]
	Cho hàm số $f(x)$ xác định trên $\mathbb{R}$, biết $f(x)=\sqrt{x^2}$. Giá trị  $f(0)$ bằng
	 \choice
		{\True $0$} 
		{$2$} 
		{$1$} 
		{không tồn tại} 
	 \loigiai{
	 Ta có $f(0)=\sqrt{0^2}=0$.}
\end{ex}

% nguồn: MaHoaNH_TracNghiem_Lop10
\begin{ex}%[0D3N1-3]%[ma:68d81a1d]
	 Cho hàm số $y=f(x)=\heva{&\dfrac{2}{x-1}, &x\in(-\infty;2]\\&x^2-1, &x\in(2;5]}$. Tính $f(3)$.
	 \choice
		{\True $8$} 
		 {$1$}
		{$7$} 
		 {$2$}
	\loigiai{ 
		Với $x=3\in (2;5]$ thì hàm số có dạng $f(x)=x^2-1$.\\ 
	Khi đó  $f(3)=3^2-1=8$.}
\end{ex}

% nguồn: MaHoaNH_TracNghiem_Lop10
\begin{ex}%[0D3N1-3]%[ma:f09216c5]
	Hàm số $f(x)=\heva{&\dfrac{2\sqrt{x-2}-3}{x-1}, &\text{~khi~}  \, x\ge 2\\&x^2+2, &\text{~khi~} \, x < 2}$. Tính $P=f(2)+f(-2)$.
	 \choice
		 {\True $P=3$}
		 {$P=\dfrac{7}{3}$}
		{$P=6$} 
		{$P=2$} 
	\loigiai{ 
		Khi $x=2$  thì $f(x)=\dfrac{2\sqrt{x-2}-3}{x-1}$.\\
		Khi  $x=-2$ thì $f(x)=x^2+2$.\\
	Từ đó  ta có $P=f(2)+f(-2)=\dfrac{2\sqrt{2-2}-3}{2-1}+(-2)^2+2=3$.}
\end{ex}

% nguồn: MaHoaNH_TracNghiem_Lop10
\begin{ex}%[0D3Y1-3]%[ma:710bad17]
	Xét một phép thử có không gian mẫu $\Omega $ và $A$ là một biến cố của phép thử đó. Phát biểu  nào sau đây \textbf{sai}?
	 \choice
		{$0\le P\left(A\right)\le 1$} 
		{Xác suất  của biến cố $A$ là $P\left(A\right)=\dfrac{n\left(A\right)}{n\left(\Omega\right)}$}
		 {$P\left(A\right)=1-P\left(\bar{A}\right)$}
		{\True $P\left(A\right)>1$} 
	 \loigiai{ Vì $0\le P\left(A\right)\le 1$ nên $P\left(A\right)>1$ là sai
	 }
\end{ex}

% nguồn: MaHoaNH_TracNghiem_Lop10
\begin{ex}%[0D3Y1-3]%[ma:5d055761]
	Cho hàm số  $y=f(x)=2023x-2022$. Tính $f(1)$.
	 \choice
		 {$f(1)=2$}
		 {$f(1)=2023$}
		{$f(1)=2022$} 
		 {\True $f(1)=1$}
	\loigiai{ 
		Ta có:  $f(1)=2023 \cdot 1-2022=1$.
	} 
\end{ex}


%%%%%%%%%%%%%%%% ĐÚNG/SAI %%%%%%%%%%%%%%%%

% nguồn: MaHoaNH_DungSai_Lop10
\begin{ex}%[0D3H1-3]%[ma:342b6107]
	Cho hàm số $y=-2x^2+x$. Khi đó: 
	 \choiceTF
		{Điểm  $(0;0)$ và $(2;-5)$ thuộc đồ thị hàm số đã cho}
		{Điểm $(1;-1)$ không  thuộc đồ thị hàm số đã cho}
		{\True Điểm thuộc đồ thị hàm số có hoành  độ $-1$ là $(-1;-3)$}
		{\True Những điểm thuộc đồ thị hàm số  có tung độ bằng 0 là $(0;0)$ và $\left(\dfrac{1}{2};0\right)$}
	 \loigiai{
		\begin{itemchoice}
			\itemch Sai. Điểm $(0;0)$ thuộc đồ  thị hàm số vì $0=-2 \cdot 0^2+0$ (đúng).\\
			Điểm $(2;-5)$ không thuộc đồ thị hàm số vì  $-5=-2.2^2+2$ (sai).
			\itemch Sai. Điểm $(1;-1)$ thuộc đồ thị hàm số vì $-1=-2\cdot{1^2}+1$  (đúng).
			\itemch Đúng. Với $x=-1$ thì $y=-2  \cdot 1^2-1=-3$, ta có điểm $(-1;-3)$ thuộc đồ thị.
			\itemch Đúng. Với $y=0$ thì $0=-2x^2+x\Leftrightarrow x(-2x+1)=0 \Leftrightarrow \hoac{&x=0\\ &-2x+1=0}  \Leftrightarrow \hoac{&x=0\\ &x=\dfrac{1}{2}.}$\\
			Vậy có hai điểm thuộc đồ thị có tung độ bằng $0$ là  $(0;0)$ và $\left(\dfrac{1}{2};0\right)$.
	\end{itemchoice}}
\end{ex}


%%%%%%%%%%%%%%%% TRẢ LỜI NGẮN %%%%%%%%%%%%%%%%

% nguồn: MaHoaNH_TraLoiNgan_Lop10
\begin{bt}%[0D3N1-3]%[ma:35543cd9]
	Giá trị nhỏ nhất của biểu thức $F\left(x; y\right)=x-2y$ với điều kiện  $\heva{&0 \leq y \leq 5\\&x \leq 0\\&x+y-2\geq 0\\&x-y-2\leq 0} $ là
	 \shortans{-10}
	\loigiai{ 
		Vẽ các đường thẳng $d_1: y=5$, $d_2: x+y-2=0$, $d_3: x-y-2=0$, $Ox: y=0$,  $Oy: x=0$.
		\begin{center}
			\begin{tikzpicture}[scale=0.6,thick,>=stealth']
					\draw[->] (-2,0)  -- (9,0) node[above]{$x$};
					\foreach \x in {1,2,3,4,5,6,7,8} 
					\draw[shift={(\x,0)},color=black] (0pt,2pt) --  (0pt,-2pt) node[below] {\footnotesize $\x$};
					\draw[->,color=black] (0,-3) --  (0,7) node[right]{$y$};
					 \foreach \y in {1,2,3,4,5,6,7}
					\draw[shift={(0,\y)},color=black] (2pt,0pt) -- (-2pt,0pt) node[left]  {\footnotesize $\y$};
					 \node[below left] at (0,0) {$O$};
					 \clip (-2,-3) rectangle (9,7);
					\draw[line width=1.2pt,smooth,samples=100,domain=-1:9] plot  (\x,{-(\x)+2});
					\draw[line width=1.2pt,smooth,samples=100,domain=-1:15.3] plot (\x,{(\x)-2}); 
					\draw[line width=1.2pt,smooth,samples=100,domain=-2:8] plot  (\x,{5});
					\fill[pattern=north east lines]  (0,2)--(0,5) --(7,5)--(2,0)-- cycle;
					\node[below right]  at (0,5) {$A$};
					\node[right] at (0,2)  {$B$};
					\node[above] at  (2,0) {$C$};
					\node[below right] at (7,5) {$D$}; 
				\end{tikzpicture}
		\end{center}
		Các đường thẳng trên đôi một cắt nhau tại  $A\left(0; 5\right)$, $B\left(0; 2\right)$, $C\left(2; 0\right)$, $D\left(7; 5\right)$.\\
		Vì điểm $M_0\left(2;  1\right)$ có toạ độ thoả mãn tất cả các bất phương trình trong hệ nên ta tô đậm các nửa mặt phẳng bờ $d_1$, $d_2$, $d_3$, $Ox$, $Oy$ chứa điểm $M_0\left(2; 1\right)$. Miền bị tô đậm là đa giác $ABCD$ kể cả các cạnh (hình bên) là miền nghiệm của hệ phương trình đã cho.\\
		Kí hiệu $F\left(A\right)=F\left(x_A; y_A\right)=x_A-2y_A$.\\ 
		 Ta có
		$F\left(A\right)=-10$, $F\left(B\right)=-4$, $F\left(C\right)=2$, $F\left(D\right)=-3$.\\ 
		 Vậy giá trị nhỏ nhất cần tìm là $F\left(A\right)=-10$.
	 }
\end{bt}

% nguồn: MaHoaNH_TraLoiNgan_Lop10
\begin{bt}%[0D3N1-3]%[ma:40ad8784]
	Biết miền nghiệm của hệ bất phương trình $\heva{&3x+y\leq 6\\&x+y\leq 4\\&x,\,y\geq 0}$ là miền tứ giác  $OABC$ như hình vẽ dưới đây
	\begin{center}
		\begin{tikzpicture}[scale=0.6,thick,>=stealth']
				\draw[->] (-3,0) --  (7,0) node[above]{$x$};
				\foreach \x in {1,2,3,4,5,6,} 
				\draw[shift={(\x,0)},color=black] (0pt,2pt) --  (0pt,-2pt) node[below] {\footnotesize $\x$};
				 \draw[->,color=black] (0,-2) -- (0,8) node[right]{$y$};
				\foreach  \y in {1,2,3,4,5,6,7}
				\draw[shift={(0,\y)},color=black] (2pt,0pt) -- (-2pt,0pt) node[left]  {\footnotesize $\y$};
				\node[below left] at (0,0)  {$O$};
				\clip (-3,-2)  rectangle (12,8);
				 \draw[line width=1.2pt,smooth,samples=100,domain=-1:9] plot (\x,{-3*(\x)+6});
				\draw[line width=1.2pt,smooth,samples=100,domain=-1:15.3] plot (\x,{-(\x)+4}); 
				\fill[pattern=north east  lines] (0,0)--(0,4) --(1,3)--(2,0)-- cycle;
				\node[above] at  (0,4) {$A$};
				\node[right] at  (1,3) {$B$};
				\node[above] at (2,0) {$C$}; 
			\end{tikzpicture}
	\end{center}
	Gọi $\left(x_0;y_0\right)$ là một nghiệm của hệ  bất phương trình trên sao cho biểu thức $P=2x+y$ đạt giá trị lớn nhất. Khi đó tổng $T=3x_0+7y_0$ bằng bao nhiêu?
	 \shortans{24}
	 \loigiai{
		Ta có: $O\left(0;0\right)$, $A\left(0;4\right)$, $B\left(1;3\right)$, $C\left(2;0\right)$. \\ 
		Biểu thức $P=2x+y$ đạt giá trị lớn nhất tại một trong các đỉnh của tứ giác  $OABC$.\\
		Tại $O\left(0;0\right)$, ta có $P=2x+y=0$.  \\
		Tại $A\left(0;4\right)$, ta  có $P=2x+y=4$. \\
		Tại $B\left(1;3\right)$,  ta có $P=2x+y=5$.\\
		Tại $C\left(2;0\right)$, ta có  $P=2x+y=4$.\\
		Ta thấy biểu thức $P=2x+y$ đạt  giá trị lớn nhất tại $B\left(1;3\right)$.\\
		Vậy $T=3x_0+7y_0=3\cdot1+7\cdot3=24$. 
	 }
\end{bt}

% nguồn: MaHoaNH_TraLoiNgan_Lop10
\begin{bt}%[0D3N1-3]%[ma:64612ca4]
	Biểu thức $L=y-x$, với $x$ và $y$ thỏa mãn hệ bất phương trình $\heva{&2x+3y-6 \leq 0\\&x\geq 0\\&2x-3y-1\leq 0}$, đạt giá trị lớn  nhất
	là $a$ và đạt giá trị nhỏ nhất là $b$. Khi đó $a+12b$ bằng  bao nhiêu?
	\shortans{-9} 
	\loigiai{ 
		Trước hết, ta vẽ ba  đường thẳng $d_1: 2x+3y-6=0$, $d_2:x=0$, $d_3: 2x-3y-1=0$.
		\begin{center}
			\begin{tikzpicture}[line join=round,line cap=round,>=stealth,font=\footnotesize,scale=0.6]
					\draw (0,0) node [above  left] {$O$};
					\draw [->] (-3.5,0)--(3.5,0) node[below]{$x$}; 
					\draw[->] (0,-3.5)--(0,3.5) node[left]{$y$}; 
					\clip (-3.5,-3.5) rectangle  (3.5,3.5);
					\draw[samples=100,smooth,domain=-3.5:3.5,black] plot  (\x,{-(2/3)*(\x)+2});
					 \draw[samples=100,smooth,domain=-3.5:3.5,black] plot (\x,{(2/3)*(\x)-(1/3)});
					 \fill[pattern=north east lines,opacity=0.6] plot[domain=-3.5:3.5](\x,{-(2/3)*(\x)+2})--plot[domain=3.5:-3.5](\x,{3.5})--cycle;
					\fill[pattern=north  east lines,opacity=0.6] plot[domain=-3.5:3.5](\x,{(2/3)*(\x)-(1/3)})--plot[domain=3.5:-3.5](\x,{-3.5})--cycle;
					\fill[pattern=north  east lines,opacity=0.6] (-3.5,-3.5)--(0,-3.5)--(0,3.5)--(-3.5,3.5)--cycle;
					\node[right] at (0,2) {$A$}; 
					\path  (0,2) node[left]{$2$};
					\node[below] at (7/4,5/6) {$B$}; 
					\draw[dashed](7/4,0) node[below]{$\tfrac{7}{4}$; 
						$\tfrac{7}{4}$}--(7/4,5/6)--(0,5/6)  node[left]{$\tfrac{5}{6}$};
					\node[right]  at (0,-1/3) {$C$};
					\path  (0,-1/3) node[left]{$\tfrac{-1}{3}$};
				\end{tikzpicture}
		\end{center}
		Ta thấy $\left(0;0\right)$ là nghiệm của cả ba bất phương trình. Điều đó có nghĩa gốc tọa độ thuộc cả ba miền nghiệm của cả ba bất phương trình. Sau khi gạch bỏ các miền không thích hợp, miền không bị gạch là miền nghiệm  của hệ (kể cả biên).
		Miền nghiệm là hình  tam giác $ABC$ (kể cả biên), với $A\left(0;2\right)$, $B\left(\dfrac{7}{4};\dfrac{5}{6}\right)$, $C\left(0;-\dfrac{1}{3}\right)$.\\
		Ta có $L\left(A\right)=L\left(2;0\right)=2-0=2$,\\  $L\left(B\right)=L\left(\dfrac{7}{4};\dfrac{5}{6} \right)=\dfrac{6}{5}-\dfrac{7}{4}=-\dfrac{11}{12}$,\\ $L\left(C\right)=L\left(0;-\dfrac{1}{3} \right)=-\dfrac{1}{3}-0=-\dfrac{1}{3}$.\\
		Vậy $a=2$,  $b=-\dfrac{11}{12}$.\\
		 Khi đó $a+12b=-9$.
	 }
\end{bt}

% nguồn: MaHoaNH_TraLoiNgan_Lop10
\begin{bt}%[0D3N1-3]%[ma:14bd9964]
	Cho các giá trị $x$, $y$ thỏa mãn điều kiện $\heva{&x-y+2 \geq 0\\&2x-y-1\leq 0\\&3x-y-2\geq 0}$. Tìm giá  trị lớn nhất của biểu thức $T=3x+2y$.
	 \shortans{19}
	\loigiai{ 
		\begin{center}
			\begin{tikzpicture}[line join=round,line cap=round,>=stealth,font=\footnotesize,scale=0.6,yscale=0.6]
					\draw  [pattern=north west lines] (-4,-5)--(-4,10)--(5,10)--(5,-5)--cycle;
					\fill[white] (1,1)--(2,4) --(3,5)-- cycle; 
					\draw (0,0) node [above left] {$O$}; 
					\draw [->]  (-4,0)--(5,0) node[below]{$x$};
					\draw[->] (0,-5)--(0,10) node[left]{$y$}; 
					\clip (-4,-5) rectangle  (5,10);
					\draw[samples=100,smooth,domain=-4:5,black]  plot (\x,{(\x)+2});
					\draw[samples=100,smooth,domain=-4:5,black] plot (\x,{2*(\x)-1}); 
					\draw[samples=100,smooth,domain=-4:5,black] plot (\x,{3*(\x)-2}); 
					\path  (1,0) node[below]{$1$}
						(0,1) node[left]{$1$} 
						(2,0)  node[below]{$2$}
						(0,4)  node[left]
						 {$4$}
						 (3,0) node[below]{$3$}
						 (0,5) node[left]{$5$};
					\node[right] at  (1,1) {$A$};
					\node[left]  at (2,4) {$B$};
					\node[right] at  (3,5) {$C$};
				\end{tikzpicture}
		\end{center}
		Miền nghiệm của hệ đã cho là miền trong tam giác $ABC$ (Kể cả đường biên) trong đó  $A\left(1;1\right)$, $B\left(2;4\right)$, $C\left(3;5\right)$.
		Giá trị lớn nhất  của $T=3x+2y$ đạt được tại các đỉnh của tam giác $ABC$.\\
		Ta có $T\left(A\right)=T\left(1;1\right)=5$, $T\left(B\right)=T\left(2;4 \right)=14$,  $T\left(C\right)=T\left(3;5\right)=19$.\\
		Vậy giá trị  lớn nhất của biểu thức $T=3x+2y=19$.
	 }
\end{bt}

% nguồn: MaHoaNH_TraLoiNgan_Lop10
\begin{bt}%[0D3N1-3]%[ma:a44bd27a]
	Gia đình chị Minh dự định trồng rau và hoa trên một mảnh đất có diện tích $8$ ha. Nếu trồng $1$ ha rau thì cần $20$ ngày công và thu lợi $3$ triệu đồng. Nếu trồng $1$ ha hoa thì cần $30$ ngày công và thu lợi $4$ triệu đồng. Biết rằng, gia đình chị Minh chỉ có thể sử dụng không quá $180$ ngày công cho công việc trồng rau và hoa. Hỏi từ việc trồng rau và hoa nói trên, chị Minh có thể thu  về lợi nhuận cao nhất là bao nhiêu triệu đồng?
	\shortans{26} 
	 \loigiai{
		Gọi $x, \, y\, \left(x\geq0, \,y\geq0\right)$ lần lượt là số  ha đất trồng rau và hoa.\\
		Diện tích đất trồng canh tác không vượt quá $8$ ha  nên ta có $x+y\leq8$.\\
		Số ngày công sử dụng không vượt quá $180$ ngày  nên $2x+3y\leq180$.\\
		Từ đó, ta  có hệ bất phương trình $\heva{&x+y\leq8\\&2x+3y\leq180\\&x\geq 0\\&y\geq 0.}$\\
		Ta cần tìm $x, \, y$ sao cho $T\left(x,y\right)=3x+4y$  lớn nhất.\\
		Miền nghiệm của hệ được biểu diễn  như sau:
		\begin{center}
			\begin{tikzpicture}[scale=0.6,thick,>=stealth']
					\draw[->] (-3,0) -- (12,0) node[above]{$x$}; 
					\foreach \x  in {1,2,3,4,5,6,7,8,9,10,11}
					\draw[shift={(\x,0)},color=black] (0pt,2pt) -- (0pt,-2pt) node[below] {\footnotesize  $\x$};
					 \draw[->,color=black] (0,-2) -- (0,9) node[right]{$y$};
					\foreach \y in  {1,2,3,4,5,6,7,8,9}
					\draw[shift={(0,\y)},color=black] (2pt,0pt) --  (-2pt,0pt) node[left] {\footnotesize $\y$};
					\node[below left] at (0,0)  {$O$};
					\clip (-3,-2)  rectangle (12,9);
					\draw[line width=1.2pt,smooth,samples=100,domain=-3:9]  plot (\x,{-(\x)+8});
					\draw[line width=1.2pt,smooth,samples=100,domain=-1:12]  plot (\x,{-(2/3)*(\x)+6});
					\fill[pattern=north  east lines] (0,0)--(0,6) --(6,2)--(8,0)-- cycle;
					\node[above] at (8,0)  {$A$};
					\node[right] at (6,2) {$B$}; 
					\node[right] at (0,6)  {$C$};
				\end{tikzpicture}
		\end{center}
		Miền nghiệm của hệ bất phương trình trên là miền trong của tứ giác $OABC$, kể cả các cạnh của tứ giác đó,  với $O\left(0;0\right)$, $A\left(8;0\right)$, $B\left(6;2\right)$, $C\left(0;6\right)$. \\
		 Tại $O\left(0;0\right)$, ta có $T=3\cdot0+4\cdot0=0$.\\
		 Tại $A\left(8;0\right)$, ta có $T=3\cdot8+4\cdot0=24$.\\
		Tại $B\left(6;2\right)$, ta có  $T=3\cdot6+4\cdot2=26$.\\
		Tại $C\left(0;6\right)$,  ta có $T=3\cdot0+4\cdot6=24$.\\
		Vậy số lợi nhuận cao nhất mà gia đình chị Minh thu được từ trồng rau và hoa là  $26$ triệu đồng.
	 }
\end{bt}

% nguồn: MaHoaNH_TraLoiNgan_Lop10
\begin{bt}%[0D3N1-3]%[ma:e4a89d6f]
	Một công ty điện tử sản xuất hai kiểu radio trên hai dây chuyền độc lập. Radio kiểu một sản xuất trên dây chuyền một với công suất $45$ radio/ngày,  radio kiểu hai sản xuất trên dây chuyền hai với công suất $80$ radio/ngày. Để sản xuất một chiếc radio kiểu một cần $12$ linh kiện, để sản xuất một chiếc radio kiểu hai cần $9$ linh kiện. Tiền lãi khi bán một chiếc radio kiểu một là $250\,000$ đồng, lãi thu được khi bán một chiếc radio kiểu hai là $180\,000$ đồng. Biết rằng số linh kiện có thể sử dụng tối đa trong một ngày là $900$. Gọi $x_0$, $y_0$ lần lượt là số radio kiểu một và radio kiểu hai sản suất được trong một ngày để tiền lãi thu được là nhiều nhất. Khi đó tổng $T=x_0+2y_0$ bằng
	 \shortans{125}
	\loigiai{ 
		Gọi $x$  và $y$ lần lượt là số radio kiểu một và số radio kiểu hai mà công ty này sản xuất trong một ngày $\left(x,\, y \in \mathbb{N}^*\right)$.\\
		Số tiền  lãi mà công ty này thu về hàng ngày là $P=250\,000x+180\,000y$ (đồng).\\
		Ta có hệ bất phương trình $\heva{&12x+9y\leq 900\\&0\leq x\leq 45\\&0\leq y\leq  80.}$\\
		Bài toán trở thành tìm giá trị lớn nhất của biểu thức $P=250\,000x+180\,000y$  trên miền nghiệm của hệ bất phương trình $\heva{&12x+9y\leq 900\\&0\leq x\leq 45\\&0\leq y\leq 80.}$
		\begin{center}
			\begin{tikzpicture}[scale=0.6,thick,>=stealth']
					 \draw[->] (-3,0) -- (10,0) node[above]{$x$};
					 \draw[->,color=black] (0,-2) -- (0,11) node[right]{$y$};
					\node[below left] at (0,0)  {$O$};
					 \clip (-3,-2) rectangle (10,11);
					\draw[line width=1.2pt,smooth,samples=100,domain=-1:9]  plot (\x,{-(12/9)*(\x)+10});
					\draw[line width=1.2pt,smooth,samples=100,domain=-1:15.3] plot  (\x,{8});
					\fill[pattern=north east lines] (0,0)--(0,8) --(1.5,8)--(4.5,4)--  (4.5,0)-- cycle;
					\node[above right] at  (4.5,0) {$A$};
					\node[below right] at (4.5,0) {$45$}; \node[below] at (7.5,0) {$75$}; 
					\node[right] at (4.5,4) {$B$}; \node[below] at  (1.5,0) {$15$};
					 \node[below left] at (0,8) {$80$};
					\node[above right] at (1.5,8) {$C$}; \node[left] at (0,4) {$40$}; 
					\node[left] at  (0,10) {$100$};
					\node[above left] at  (0,8) {$D$};
					 \draw[dashed] (0,4)--(4.5,4)
						 (1.5,0)--(1.5,8);
					\draw[thick]  (4.5,-2)--(4.5,11);
				\end{tikzpicture}
		\end{center}
		Miền nghiệm của hệ bất phương trình $\heva{&12x+9y\leq 900\\&0\leq x\leq 45\\&0\leq y\leq 80}$ là miền ngũ giác  $OABCD$ trong đó $O\left(0;0\right)$, $A\left(45;0\right)$, $B\left(45;40\right)$, $C\left(15;80\right)$, $D\left(0;80\right)$. \\
		Tại $O\left(0;0\right)$  , ta có $P=250\,000\cdot0+180\,000\cdot0=0$.\\
		Tại $A\left(45;0\right)$ ta  có $P=250\,000\cdot45+180\,000\cdot0=11\,250\,000$.\\
		Tại $B\left(45;40\right)$ ta  có $P=250\,000\cdot45+180\,000\cdot40=18\,450\,000$. \\
		Tại $C\left(15;80\right)$ ta có  $P=250\,000\cdot15+180\,000\cdot80=18\,150\,000$. \\
		Tại $D\left(0;80\right)$ ta có $P=250\,000\cdot0+180\,000\cdot80=14\,400\,000$.\\ 
		 Ta thấy biểu thức $P=250\,000x+180\,000y$ đạt giác trị lớn nhất khi $x_0=45, \, y_0=45$.\\
		Vậy  $T=x_0+2y_0=125$.
	 }
\end{bt}

% nguồn: MaHoaNH_TraLoiNgan_Lop10
\begin{bt}%[0D3H1-3]%[ma:4bffd3df]
	Cho hàm số $y=f(x)=m{x^3}-2(m^2+1){x^2}+2m^2-m$ với $m$ là tham số. Tìm tổng tất cả các giá trị của $m$ sao cho  $f(-1)=2$.
	\shortans{-2} 
	\loigiai{ 
		Ta có $f(-1)=2 \Leftrightarrow 2=-m-2(m^2+1)+2m^2-m  \Leftrightarrow m=-2$.\\
	Vậy tổng tất cả các giá trị  của $m$ bằng $-2$.}
\end{bt}

% nguồn: MaHoaNH_TraLoiNgan_Lop10
\begin{bt}%[0D3H1-3]%[ma:6a621b1c]
	Cho $f(x)=\heva{&x-4, &\text{~khi~}  \, x\ge 0\\ &x^2-4x+1, &\text{~khi~} \, x < 0}$. Có bao nhiêu giá trị của tham số $m$ để $f(m^2)+f(-2)=18$.
	\shortans{2} 
	\loigiai{ 
		Với $x=m^2\ge  0$, ta có $f(m^2)=m^2-4$.\\
		Với $x=-2 < 0$,  ta có $f(-2)=(-2)^2-4(-2)+1=13$.\\
		Khi đó $f(m^2)+f(-2)=18 \Leftrightarrow (m^2-4)+13=18 \Leftrightarrow m=\pm  3$.\\
	Vậy có $2$ giá trị của $m$ thỏa mãn yêu cầu bài  toán.}
\end{bt}


%%%%%%%%%%%%%%%% TỰ LUẬN %%%%%%%%%%%%%%%%

% nguồn: MaHoaNH_TuLuan_Lop10
\begin{bt}%[0D3H1-3]%[ma:5cb5cb78]
	Cho hàm số $f(x)=\heva{&1, &\text{~khi~} \, x=2\\ &\dfrac{1}{x-2}, &\text{~khi~} \, x\ne 2}$.  Giải phương trình $f(x)=\dfrac{x+1}{x-2}$.
	 \loigiai{
		 Xét phương trình $f(x)=\dfrac{x+1}{x-2}$.\\
		Điều kiện $x\ne 2$,  khi đó $f(x)=\dfrac{1}{x-2}$.\\
		Ta có  phương trình $\dfrac{1}{x-2}=\dfrac{x+1}{x-2} \Leftrightarrow x+1=1\Leftrightarrow x=0$.\\
		Vậy tập nghiệm  phương trình $S=\{ 0\}$.
	 }
\end{bt}
